\documentclass[11pt]{article}
\usepackage[margin=1in]{geometry}
\usepackage{amsmath,amssymb,amsthm,enumitem,hyperref}
\hypersetup{colorlinks=true,linkcolor=blue,urlcolor=blue,citecolor=blue}
\newtheorem{theorem}{Theorem}
\newcommand{\HH}{\mathcal H}
\title{Matching crossing polynomials fail third-order $q$-log-convexity\\\large MathDB 317480}
\author{Principia Math In Progress}
\date{September 21, 2026}
\begin{document}
\maketitle
\begin{abstract}
For $F_n(x)=M_{2n}(x)$, the generating polynomial for the
maximum number of mutually crossing arcs in a perfect matching,
the third centered log-convexity iterate at index $5$ has degree
$34$ and leading coefficient $-14224$. This disproves Chen's
Conjecture 1.8. The input polynomials are already tabulated in
the source and are independently verified by exhaustive enumeration.
\end{abstract}

\section*{Task 1. Full disproof}
A perfect matching on the ordered set $[2n]$ has crossing
number $k$ if $k$ is the maximum number of mutually crossing
arcs in its standard arc diagram. In particular, two arcs
$(i,j)$ and $(a,b)$ with $i<j$ and $a<b$ cross precisely when
$i<a<j<b$ or $a<i<b<j$. Let
\[
 F_n(x)=M_{2n}(x)=\sum_k M_{2n,k}x^k.
\]
We apply the source's centered operator across $n$:
\[
 (\HH A)_n=A_{n-1}A_{n+1}-A_n^2.
\]
Conjecture 1.8 of \cite{Chen} asserts that every iterate is
coefficientwise nonnegative, and proposes a further
log-concavity property for those nonnegative iterates.

\begin{theorem}
The matching polynomials satisfy
\[
 (\HH^3F)_5(x)=-14224x^{34}
              +\text{a polynomial of degree at most }33.       \tag{1}
\]
They are therefore not infinitely $q$-log-convex.
\end{theorem}

\begin{proof}
The initial polynomials printed on p.~6 of \cite{Chen} give
\[
 F_n(x)=x^n+a_nx^{n-1}+b_nx^{n-2}
        +\text{lower-degree terms},
\]
with
\[
\begin{array}{c|rrrrrrr}
n&2&3&4&5&6&7&8\\\hline
a_n&2&9&20&35&54&77&104\\
b_n&0&5&70&315&891&2002&3900
\end{array}.                                                    \tag{2}
\]
The unique matching with all $n$ arcs mutually crossing has
pairs $(1,n+1),\ldots,(n,2n)$, accounting for the leading $1$.
The remaining entries in (2) are known source data; the complete
independent enumeration certificate is described below.

Put $B_n=(\HH F)_n$ and $C_n=(\HH^2F)_n$.
The leading terms of degree $2n$ cancel in $B_n$. For
$3\le n\le7$, table (2) gives
\[
 a_{n-1}+a_{n+1}-2a_n=4,
\]
so
\[
 B_n(x)=4x^{2n-1}+v_nx^{2n-2}+\text{lower-degree terms},
\]
where
\[
 v_n=b_{n-1}+b_{n+1}+a_{n-1}a_{n+1}-a_n^2-2b_n.
\]
Substitution yields
\[
\begin{array}{c|rrrrr}
n&3&4&5&6&7\\\hline
v_n&19&95&186&314&474
\end{array}.                                                    \tag{3}
\]
In $C_n=B_{n-1}B_{n+1}-B_n^2$, the terms of degree $4n-2$
cancel. The coefficient of the next power $x^{4n-3}$ is
\[
 c_n=4(v_{n-1}+v_{n+1}-2v_n).
\]
Hence
\[
 C_4=60x^{13}+\cdots,\qquad
 C_5=148x^{17}+\cdots,\qquad
 C_6=128x^{21}+\cdots.                                        \tag{4}
\]
It follows that $C_4C_6-C_5^2$ has degree $34$ and leading
coefficient
\[
 60\cdot128-148^2=7680-21904=-14224.
\]
This proves (1) and refutes infinite $q$-log-convexity.
\end{proof}

This is an interior counterexample: only $F_2,\ldots,F_8$,
or $M_4,\ldots,M_{16}$, are involved. At an ordinary positive
parameter it gives the stronger numerical witness
\[
 (\HH^3F)_5(10)=
 -1584481155042307779051353803951488000000000<0.                \tag{5}
\]

\section*{Independent counting certificate}
The verifier enumerates all $2173623$ matchings on the sets
$[4],[6],\ldots,[16]$. At each recursive step it takes the
smallest unmatched vertex $i$ and pairs it with every remaining
vertex $j$. Induction on the number of unmatched vertices shows
that every matching occurs once and only once.

For each matching, it forms the graph whose vertices are its
arcs and whose edges join pairs with alternating endpoints.
Its maximum clique size is exactly the crossing number by
definition. An exhaustive clique search therefore determines
the entire coefficient row, not just the three coefficients
used in (2). The row totals are $(2n-1)!!$, and every row agrees
with the source. Independent hook-length sums over partitions
with even row lengths give the same seven distributions.

Two exact polynomial multiplication methods reproduce all
three iterates, including the cancellations in (3)--(4).
Direct integer evaluation checks (5). No floating point or
asymptotic approximation enters the verification. Review is
by the authoring agent only.

The source's reported third-order positivity evidence is
incompatible with (1), although its printed polynomials agree
with the direct counts. The stronger regularity assertion on
p.~7 explicitly starts with coefficientwise positivity, which
already fails here. We do not claim that a quadratic
log-concavity inequality must fail for the resulting signed
coefficient list. Ordinary and second-order $q$-log-convexity,
and log-concavity of each original coefficient row, remain
separate questions.

\section*{Tasks 2--5. Follow-up directions}
\begin{enumerate}[start=2,leftmargin=*]
\item Decide second-order $q$-log-convexity of the matching
crossing polynomials.
\item Determine all positive parameters $x$ for which
$(M_{2n}(x))_{n\ge1}$ is infinitely log-convex.
\item Determine whether each fixed iterate is eventually
coefficientwise nonnegative as $n$ increases.
\item Derive formulas for the coefficients nearest the highest
power in $\HH^rF$, and classify the resulting sign changes.
\end{enumerate}

\begin{thebibliography}{9}
\raggedright
\bibitem{Chen} W. Y. C. Chen,
\emph{Log-concavity and $q$-Log-convexity Conjectures on the
Longest Increasing Subsequences of Permutations},
\href{https://arxiv.org/abs/0806.3392v2}{arXiv:0806.3392v2},
June 22, 2008. Conjecture 1.8 and pp.~5--7.
\bibitem{DB} MathDB, \href{https://mathdb.com/p/317480}{Problem 317480},
accessed September 21, 2026.
\end{thebibliography}
\end{document}
